\section{Full model under consideration}
\label{sec:full_model}

The driven sweet spot developed in the preceding chapter suppresses low-frequency flux-induced dephasing by canceling the first-order sensitivity of the dressed cavity frequency to fluctuations of $\omega_b$. However, the drive also opens new dephasing channels when FTT relaxation is present. This chapter quantifies these residual channels and evaluates the total cavity dephasing rate.

To this end, the system is modeled with a stochastic master equation. For a given flux-noise realization $\delta\Phi(t)$, the density matrix $\rho(t)$ obeys
\begin{align}
\dot{\rho}
&= -i\left[\,H_{\mathrm{sys}} + \delta\omega_b(t)\, b^\dagger b,\ \rho\,\right]
+ \gamma_{b\downarrow}\,\mathcal{D}[b]\rho,
\nonumber \\
H_{\mathrm{sys}}
&= H_{\mathrm{static}} + 2\Omega_0 \cos(\omega_d t)\,(b + b^\dagger).
\label{eq:Linblad_ch8}\end{align}
Here $\mathcal{D}[b]\rho = b\rho b^\dagger - \tfrac{1}{2}\{b^\dagger b, \rho\}$ describes FTT relaxation at rate $\gamma_{b\downarrow}$ (the FTT energy-relaxation rate).

Flux noise in these devices typically exhibits a $1/f$ spectrum. The corresponding power spectral density is taken as
\begin{equation}
S_{1/f}(\omega)
= \int_{-\infty}^{\infty} d\tau\, \langle \delta\Phi(\tau)\,\delta\Phi(0)\rangle\, e^{i\omega \tau}
= \frac{S_0^{\,2}}{\left|\omega/2\pi\right|},
\end{equation}
where $S_0$ sets the flux-noise amplitude.

The following analysis is carried out in the effective-drive frame to separate the roles of relaxation and dephasing. The equation of motion is first derived in the interaction picture (Sec.~\ref{sec:eom_interaction}), and then analytical (Sec.~\ref{sec:analytical_rates}) and numerical (Sec.~\ref{sec:numerical_rates}) calculations of the dephasing rates are presented for the unselective and selective regimes in turn.
